%% ch06_experimental_setup.tex  —  Chapter 6: Experimental Setup (Step 1 polish staging)
%% HSP-Agile CCF-B Report, East China Normal University, 2026
%% Ye Xujun

\chapter{Experimental Setup}
\label{ch:experiments}

This chapter fixes the measurement protocol for the Tier-1 claims:
what is compared, how prevention is scored, and which run ids are
authoritative.
Hardware, modes, metrics, and experiment tracks are stated once;
threats are previewed briefly and analysed in Chapter~\ref{ch:discussion}.

% Compact authoritative numbers for long report / conference.
% Source: artifacts/AUTHORITATIVE_NUMBERS.md + e6_paired_summary.json
\begin{table}[t]
  \centering
  \caption{Authoritative numbers card (cite these; do not mix corpora).}
  \label{tab:authoritative-numbers}
  \footnotesize
  \begin{tabular}{@{}p{0.28\linewidth}p{0.66\linewidth}@{}}
    \toprule
    \textbf{Claim / corpus} & \textbf{Numbers} \\
    \midrule
    E6 lead (C2), $K{=}3$
      & A/B/C Conf.\ 79.1\,/\,78.9\,/\,86.9\%;
        C$-$A ${+}7.7$\,pp; paired 14/4/102;
        95\% CI $[2.3,\,13.6]$; Wilcoxon $p{=}0.018$ \\
    E14 scope (not uniqueness)
      & \texttt{semantic\_ir} 75.1\% $\neq$
        \texttt{execution\_trace\_matched} 85.4\%
        (paired 5/19/96) \\
    E2 prevention (C3)
      & M PDR 95.0\% / FAR 5.0\% vs.\ B2 91.2\% / 8.8\%
        (impl-screening; bootstrap CI excludes 0) \\
    Fixed-oracle stress (bundle $K{=}5$)
      & E1: B1 84.2 / B2 98.3 / M 100.0
        (${+}1.7$\,pp; 2/120 wins);
        E10 ${+}1.0$\,pp; E12 B2 100.0 / M 99.7
        (B2 first every seed) \\
    Equal-$K$ Conf.\ (\texttt{M\_eq})
      & \texttt{run\_e1\_equal\_k\_v1}: B2 97.5\% /
        \texttt{M\_eq} 100.0\% (${+}2.5$\,pp; 3/0/117) \\
    Industrial C4
      & $n{=}31$: B2 100\% Conf.\ under gpt-4o / Claude;
        B1/M 95.5\% --- prefer B2 by default \\
    \bottomrule
  \end{tabular}
\end{table}

% Master protocol table for the CCF-B Accept revision.
% Keep this as the single protocol map for conference and full versions.
\begin{table}[t]
  \centering
  \caption{Master protocol table.
    Each headline number must trace to one row; this prevents mixing $K$,
    feedback renderers, gates, and corpora across claims.
    B6 is a Conf.-only VerifierLoop baseline on BENCH-120; prevention equals
    B2/A2 by design when Screen is off, so the DGA prevention row compares
    B2 / A2 / M.}
  \label{tab:master-protocol}
  \footnotesize
  \renewcommand{\arraystretch}{1.08}
  \begin{tabular}{@{}p{0.12\linewidth}p{0.19\linewidth}p{0.06\linewidth}p{0.19\linewidth}p{0.18\linewidth}p{0.18\linewidth}@{}}
    \toprule
    Claim & Comparison & $K$ & Feedback / gate & Corpus & Primary metric \\
    \midrule
    DGA / Decomp
      & B2 / A2 / M
      & 3--5
      & Mode defaults; gate reported explicitly
      & E2 faulty mutants
      & PDR / FAR by eval type \\
    SIFR
      & A / B / C under full M
      & 3
      & test-only / expected / Semantic IR
      & BENCH-120 and headroom strata
      & $\Delta$Conf., W/L/T, CI \\
    External scope
      & A / B / C or B2 / M\_eq
      & 3
      & Same as source run
      & RealSpec / HKCA09 / GitHub / desensitized
      & Conf. and saturation boundary \\
    HTE
      & B2 default vs. M escalation
      & 3--5
      & Pre-hoc rule, no attempt-history features
      & Pooled public + synthetic metadata
      & Escalation precision / recall \\
    Stress only
      & B1 / B2 / strengthened M
      & 3 / 5
      & M uses trace-matched feedback + advisory gate
      & Fixed-oracle BENCH-120
      & Near-ceiling Conf. context \\
    \bottomrule
  \end{tabular}
\end{table}


\begin{inlinetable}{Claim-to-experiment map (Tier-1 focus).}
  \label{tab:claim-map}
  \footnotesize
  \begin{tabular}{llp{0.42\linewidth}}
    \toprule
    Claim & Experiments & What would falsify it \\
    \midrule
    Premise (SpecIR hinge) & E3, E4, E8 & Adapters fail to transfer; SMT
      coverage equals random at equal budget \\
    C2 (structured feedback) & E5, E6, E14, combo-seed & Full IR $\not>$
      test-only under headroom; E14 uniqueness over-claimed \\
    C3 (conjunctive gate) & E1 Strict, E2 PDR/FAR, A2 & High Conf.\ with
      high FAR; gate adds nothing on prevention \\
    C4 (deploy B2 default) & E3, E10, E12, pilots & Overlap tertiles alone
      dictate M; pilots contradict B2 default \\
    \bottomrule
  \end{tabular}
\end{inlinetable}

%=============================================================================
\section{Research Questions}
\label{sec:exp:rq}
%=============================================================================

Numbering is mechanism-first (same scheme as Chapters~\ref{ch:intro}
and~\ref{ch:results}): typed feedback and overlap structure, then
prevention and deployment, with aggregate E1 ranking as an RQ5
stress-test rather than a universal-effectiveness headline.

\paragraph{RQ1 (Specification-indexed Feedback IR).}
Does Semantic Feedback IR improve repair beyond unstructured test-only
feedback?
(E6, E14 scope, hard combo-seed probe; primarily C2.)

\paragraph{RQ2 (Gap vs.\ Guard-Overlap Density; negative).}
Does the M--baseline gap grow monotonically with guard-overlap density?
(E3, E10; bounds C4---not a positive SpecIR ranking claim.)

\paragraph{RQ3 (Conjunctive Accept vs.\ Conf Alone).}
Does conjunctive $\mathrm{Accept}$ improve prevention versus mean
conformance alone?
(E2, A2 vs.\ M; primarily C3.)

\paragraph{RQ4 (Deployment Boundary).}
Under what conditions should practitioners deploy M versus B2?
(E3, E10, E11, E8c, pattern/pilot slices; C4 /
Table~\ref{tab:deployment-boundary}.)

\paragraph{RQ5 (Quality--Overhead Trade-off).}
Is the latency cost justified by prevention and release requirements?
(E1 Pareto, A2 operating point.)
Primary fixed-oracle E1 ranking (B1/B2/M from \code{run\_e1\_m\_win\_v2})
is an overlap stress-test with paired win rates, not a powered dominance
claim.

%=============================================================================
\section{Hardware and Software Environment}
\label{sec:exp:environment}
%=============================================================================

All runs used a Windows~10 (build~26220) desktop, 6-core CPU,
Python~3.12, and one \code{multiprocessing} worker per core.
Z3~4.16~\citep{demoura2008z3,barrett2018smt} via \code{z3-solver} with a
10-second per-call timeout; \code{unknown} is non-constraining, not a hard
failure.

%=============================================================================
\section{Language Model Configuration and Run Protocol}
\label{sec:exp:protocol}
%=============================================================================

\subsection{Language Model Configuration}

LLM-based modes (B1, B2, M, A1--A3) use
Qwen2.5-27B-Instruct~\citep{yang2024qwen25} (\textsf{Qwen-27B}) through an
OpenAI-compatible gateway.
Temperatures: $T_\text{gen}{=}0.7$, $T_\text{repair}{=}0.0$; token cap 4096
per call (Chapter~\ref{ch:method}).
Repair budget~$K$ is not uniform across corpora: primary fixed-oracle ranking
(E1/E10/E12 M-win) uses strengthened M with $K{=}5$ (best-effort argmax
Conf.); B1/B2 keep baseline budgets.
Mechanism and controlled comparisons (E6; historical pre-fix
matrix) use $K{=}3$.

\subsection{Corpora and Measurement Correction}

Three corpora must stay distinct (Table~\ref{tab:corpus-card};
\texttt{artifacts/AUTHORITATIVE\_NUMBERS.md}).
\emph{Once for the chapter:} a bug in others-witness generation
(others encoded as $\mathit{True}$ without $\neg$prior guards) previously
inflated false failures (${\sim}$5\% Strict on canonical E1).
Fixed first-match others witnesses enable fair evaluation; primary B1/B2/M
aggregates cite the fixed-oracle M-win runs.
Pre-fix runs (\code{run\_e1\_canonical\_v1}/
\code{run\_hard\_full\_parallel\_v1} and related) are archive only and must
never serve as primary ranking evidence.
Fixed-oracle A1--A3 live in \code{run\_e1\_ablation\_fixed\_v1}
(Table~\ref{tab:ablation-fixed}).

\begin{inlinetable}{Corpus / protocol card (which experiment cites which run).}
  \label{tab:corpus-card}
  \footnotesize
  \begin{tabular}{p{0.22\linewidth}p{0.32\linewidth}p{0.38\linewidth}}
    \toprule
    Role & Run id(s) & Used for \\
    \midrule
    Primary ranking
      (fixed others-witness oracle)
      & \code{run\_e1\_m\_win\_v2},
        \code{run\_e10\_m\_win\_v1},
        \code{run\_e12\_m\_win\_v1},
        \code{run\_e1\_ablation\_fixed\_v1},
        \textbf{\code{run\_e1\_equal\_k\_v1}}
      & E1/E10/E12 B1/B2/M Conf stress
        (M $K{=}5$ bundle);
        \textbf{equal-$K$ Conf:\ B2 vs.\ \texttt{M\_eq}
        ($K{=}3$, ${+}2.5$\,pp)};
        fixed-oracle A1--A3 \\
    Mechanism / prevention
      & \code{run\_feedback\_v2},
        \code{prevention\_full\_v1}
      & E6 typed-IR ($K{=}3$); E2 PDR/FAR \\
    Historical pre-fix
      (superseded for ranking)
      & \code{run\_e1\_canonical\_v1}/
        \code{run\_hard\_full\_parallel\_v1};
        also \code{run\_e12\_canonical\_v1},
        \code{run\_e10\_random\_v2}
      & Pre-fix A1--A3 Conf.\ ordering
        (archive), some latency figures,
        pre-fix Holm multi-mode stats;
        $K{=}3$; buggy others-witness
        oracle ($\sim$5\% Strict) \\
    \bottomrule
  \end{tabular}
\end{inlinetable}

Primary scales: E1 $120{\times}3$ modes; E10 $100$; E12
$120{\times}3$ seeds${\times}3$ modes.
Extended baselines B3--B5 use \code{run\_ccf\_b\_extended\_v1} (repeat~0).
Artefacts (prompts, responses, Z3 logs, traces, timing) are JSONL;
task/prompt RNG seed~42; repair calls at temperature~0.0 are near-deterministic.
Full reproduction: Appendix~\ref{app:reproducibility}.

%=============================================================================
\section{Benchmark Construction}
\label{sec:exp:benchmark}
%=============================================================================

Ordered FSF artefacts are rare in public repos, so the corpus is
constructed to be discriminative (baselines fail non-trivially),
formally valid (Z3-satisfiable), and diverse in orderings and boundaries.
DafnyBench~\citep{loughridge2024dafnybench} targets annotation-heavy
verification rather than FSF precedence.

\subsection{Hard Synthetic Benchmark (E1)}

Twenty base SOFL/FSF tasks from published
examples~\citep{liu2004soflbook,li2022qrs_companion} form a development
set and are excluded from primary evaluation.
The \emph{HardTaskGenerator} emits tasks with 5 integer inputs
$(a{,}b{,}c{,}d{,}e)\in[0,100]^5$, 3 outputs, and 8 ordered scenarios with
precedence-sensitive overlapping guards (targeting SOV from
Chapter~\ref{ch:formalization}).
Z3 admits only individually satisfiable scenarios; 180 candidates yield
120 accepted tasks (60 discarded).
Figure~\ref{fig:benchmark-pipeline} shows the path; Table~\ref{tab:benchmark-stats}
summarises the set (mean guard-overlap rate 1.8 scenarios per input point).
Under fixed-oracle E1 (\code{run\_e1\_m\_win\_v2}), B1 mean Conf.\ is 84.2\%;
discriminative signal sits in per-task variance
(Figure~\ref{fig:conf-distribution}) and prevention metrics rather than in
aggregate mean Conf.\ alone.

\begin{figure}[t]
  \centering
  \tikzinput{diagrams/tikz/benchmark_pipeline}
  \caption{\textbf{Key observation:} the Z3 filter keeps only overlap-heavy,
    satisfiable tasks---the regime where ordering defects actually occur.
    \textbf{How to read:} left-to-right from sampled guards through filtering
    to the final 120-task set.}
  \label{fig:benchmark-pipeline}
\end{figure}

\begin{table}[H]
  \caption{Summary statistics of the 120-task hard benchmark.}
  \label{tab:benchmark-stats}
  \centering
  \thesistable{%
    \begin{tabular}{lc}
      \toprule
      Property & Value \\
      \midrule
      Total tasks & 120 \\
      Inputs per task & 5 \\
      Outputs per task & 3 \\
      Scenarios per task & 8 \\
      Mean guard-overlap rate (scenarios/point) & 1.8 \\
      Tasks rejected by Z3 filter & 60 \\
      Tasks accepted & 120 \\
      Conformance cases per task (\code{strict\_eval}) & 64 \\
      \bottomrule
    \end{tabular}%
  }
\end{table}

\paragraph{Illustrative skeleton.}
The hard tasks scale the first-match idea of \code{compute\_tax}
(Chapter~\ref{ch:intro}): nested guards, most-specific first, defects hiding
in overlap regions.
A reverse-order evaluation yields the wrong branch on overlapping inputs---a
scenario-order violation detectable by the pattern guard and by Z3 witnesses
(e.g.\ \code{a=85, b=15, c=10, d=80, e=20}).

\subsection{External and Auxiliary Corpora}
\label{sec:exp:e10}

\begin{inlinetable}{Auxiliary benchmarks (external validity and controls).}
  \label{tab:aux-benchmarks}
  \footnotesize
  \begin{tabular}{llp{0.48\linewidth}}
    \toprule
    Slice & $n$ / run & Role \\
    \midrule
    E10 random
      & 100; \code{run\_e10\_m\_win\_v1}
      & Same generator \emph{without} Z3 overlap filter;
        negative control for filter bias \\
    E11 external SOFL
      & 22; \code{run\_e11\_external\_v1}
      & Hand-authored from published Agile-SOFL /
        IET patterns~\citep{liu2004soflbook,li2022qrs_companion,li2023iet_faultprevention};
        $K{=}3$; IDs disjoint from E1 \\
    Real-priority micro
      & 30; \code{run\_real\_priority\_micro\_v1}
      & Non-synthetic ordered first-match (industrial, GitHub
        \texttt{.asfl}, HKCA09, desensitized pilots, adapted ACL/tax);
        equal-$K$ B1/B2/\texttt{M\_eq}; Table~\ref{tab:real-priority-micro} \\
    E8 real-derived
      & 40 (20 HumanEval~\citep{chen2021codex}{+}20 MBPP~\citep{austin2021mbpp})
      & Ordered-branch problems converted to FSF; E8c full B1/B2/M \\
    SMT domain ablation
      & 30 real-priority tasks
      & Box sizes $[-5,20]$--$[-1000,1000]$ plus Real encoding;
        Table~\ref{tab:smt-scalability},
        Figure~\ref{fig:smt-domain-latency} \\
    \bottomrule
  \end{tabular}
\end{inlinetable}

\label{sec:exp:e11}\label{sec:exp:real-priority}\label{sec:exp:smt-scalability}%
E10 asks whether aggregate M${>}$B2 holds without overlap filtering, or
whether tertiles (E3) and deployment signals
(Table~\ref{tab:deployment-boundary}) must be read jointly.
Real-priority and SMT protocols are detailed in
Sections~\ref{sec:results:real-priority}--\ref{sec:results:smt-scalability};
selection protocol for E8 is in the artifact package.

%=============================================================================
\section{Compared Methods}
\label{sec:exp:methods}
%=============================================================================

The design space spans formal reference (B0), LLM baselines (B1--B5),
verifier-in-the-loop (B6), full HSP-Agile (M), and ablations (A1--A3).
Primary fixed-oracle ranking reports B1/B2/M from \code{run\_e1\_m\_win\_v2};
historical seven-mode and B3--B5 runs supply extended baselines
(Table~\ref{tab:corpus-card}).

\begin{table}[H]
  \caption{Compared modes.
    Columns indicate whether each component is active.
    B3--B5 are evaluated in E1-ext (\code{run\_ccf\_b\_extended\_v1},
    repeat~0 merged into Table~\ref{tab:main-results}).}
  \label{tab:modes}
  \centering
  \thesistable{%
    \footnotesize
    \renewcommand{\arraystretch}{1.15}
    \begin{tabular*}{\linewidth}{@{\extracolsep{\fill}}l l c c c@{}}
      \toprule
      Mode & Description & Formal checker & Pattern guard & Repair loop \\
      \midrule
      B0 & Reference (template) & eval only & --- & --- \\
      B1 & One-shot LLM & --- & --- & --- \\
      B2 & Test-feedback & --- & --- & \checkmark (tests) \\
      B3 & Self-Refine~\citep{madaan2023selfrefine} & --- & --- & \checkmark (self-critique) \\
      B4 & Self-Debug~\citep{chen2024selfdebug} & --- & --- & \checkmark (execution trace) \\
      B5 & Reflexion-lite~\citep{shinn2023reflexion} & --- & --- & \checkmark (verbal memory) \\
      B6 & VerifierLoop-FSF (REV-7) & \checkmark & --- & \checkmark (witness cex) \\
      M  & Full \textsc{SgDP}/HSP-Agile & \checkmark & \checkmark & \checkmark (IR) \\
      A1 & Ablation: no formal checker & --- & \checkmark & \checkmark (IR) \\
      A2 & Ablation: no pattern guard & \checkmark & --- & \checkmark (IR) \\
      A3 & Ablation: no repair loop & \checkmark & \checkmark & --- \\
      \bottomrule
    \end{tabular*}%
  }
\end{table}

\paragraph{Baselines.}
B0 is a deterministic template renderer (formal upper bound; Conf.\ computed
ex post).
B1 is one-shot LLM with no loop.
B2 adds iterative FSF-derived unit-test feedback~\citep{zhu1997specification}
at $K{=}3$ (even when M uses $K{=}5$ on M-win harnesses); it has no Z3 checker and no pattern
guard.
B3--B5 isolate unstructured self-improvement
(Self-Refine~\citep{madaan2023selfrefine}, Self-Debug~\citep{chen2024selfdebug},
Reflexion-lite~\citep{shinn2023reflexion}) under equal $K{=}3$ and the same
token cap; results in Table~\ref{tab:main-results}.
B6~\citep{sun2024clover} exposes structured SMT counterexamples in the repair
prompt without Semantic Feedback IR fields or the pattern guard
(\code{run\_b6\_full\_v1}, $n{=}100$).

\paragraph{M and ablations.}
M runs the full pipeline (Chapters~\ref{ch:method}--\ref{ch:implementation}):
Z3 witnesses, pattern screen, and CEGIS repair with Semantic Feedback IR.
Primary M-win uses $K{=}5$, \texttt{execution\_trace\_matched}, and an advisory
critical-only pattern gate; E6 keeps $K{=}3$ under full M to isolate feedback
content.
A1 replaces Z3 cases with B2 tests; A2 disables the pattern guard; A3 applies
checker and guard once without repair feedback.

%=============================================================================
\section{Mutant Design}
\label{sec:exp:mutants}
%=============================================================================

Prevention metrics need faults that look locally plausible yet violate global
precedence.
\emph{Specification-confusion} mutants perturb the task (SNO order swap/negate,
ORO output reverse, MCO drop condition, BCO boundary offset).
\emph{Implementation-screening} mutants perturb the candidate (ICO flip
condition, WRO reorder returns, MBO drop branch, DRO reverse dependency).
Each applies one first-order operator to a known-correct artefact;
non-compilable mutants are excluded.
These operators model recurring LLM guard-generation failures: lost
preconditions, reversed comparisons, missing branches, boundary shifts, and
incorrect first-match order.
They are defined over guard semantics and postconditions rather than copied
from the RF01--RF14 detector catalogue.
Consequently, a mutant can evade the configured static screen yet still fail a
SpecIR witness, or trigger the screen without being tied to one hand-written
pattern.
Per-operator PDR/FAR reports this distinction instead of treating
all injected faults as one detector-friendly pool.

\paragraph{Miniature example (PDR/FAR counting).}
An SNO mutant swaps $S_1$/$S_2$ in the oracle; $\mathrm{Accept}$ rejecting a
candidate that still implements the original order is a true prevention
(PDR numerator).
A WRO mutant reorders \texttt{if} branches in the candidate; again
$\mathrm{Accept}$ should reject.
False acceptance (FAR) occurs only if a faulty mutant still passes every
witness and the pattern screen---bounded by Theorem~\ref{thm:far-bound}.

%=============================================================================
\section{Evaluation Metrics}
\label{sec:exp:metrics}
%=============================================================================

\paragraph{Conformance and strict success.}
$\mathit{Conf}(c,t)\in[0,1]$ (Definition~\ref{def:conformance}) is the fraction
of strict evaluation cases satisfied; mean $\overline{\mathit{Conf}}$ is
reported over 120 tasks.
A \emph{strict success} requires all 64 \code{strict\_eval} cases;
modes with the pattern screen additionally require a clear high-severity
screen for $\mathrm{Accept}{=}1$ (Equation~\eqref{eq:accept}).

\paragraph{Prevention and cost.}
PDR / FAR (Definitions~\ref{def:pdr}--\ref{def:far}) are fractions of
\emph{faulty} mutants correctly blocked / falsely accepted;
on a fixed faulty set, $\mathrm{PDR}+\mathrm{FAR}=1$.
Mutation kill rate $\mu_k$ measures checker precision on syntactic mutants.
LLM call count (1.0 for B1; up to~$K$ for repair modes) and wall-clock
latency include LLM, Z3, pattern guard, and orchestration.

%=============================================================================
\section{Evaluation Protocol}
\label{sec:exp:eval-protocol}
%=============================================================================

Each experiment answers one decision question; artifact paths live in
Appendix~\ref{app:reproducibility}.

\begin{inlinetable}{Experiment map (E1--E17).}
  \label{tab:eval-protocol}
  \footnotesize
  \begin{tabular}{llp{0.52\linewidth}}
    \toprule
    Exp & RQ & Decision question / protocol \\
    \midrule
    E1 & RQ5 & Aggregate stress-test: Conf/PDR/FAR/latency;
      \code{run\_e1\_m\_win\_v2} (M $K{=}5$); win rates, no Holm claim \\
    E2 & RQ3 & PDR/FAR per operator; 120${\times}$6 modes${\times}$8 ops;
      \code{prevention\_full\_v1} \\
    E3 & RQ2 & Conf by overlap tertile / guard-complexity \\
    E4 & --- & SMT vs.\ random coverage across budgets $\{4{\ldots}64\}$ \\
    E5 & RQ1 & $\mathrm{Conf}(k)$ dynamics from \code{attempt\_history} \\
    E6 & RQ1 & Feedback variants A/B/C under M, $K{=}3$;
      \code{run\_feedback\_v2} \\
    E7 & --- & Pattern-guard P/R/F1 on E2 mutants \\
    E8 & --- & SpecIR adapters: SOFL, Mini-Z, Mini-StateMachine;
      E8b $n{=}30$; E8c HumanEval/MBPP $n{=}40$ \\
    E9 & --- & Failure taxonomy of non-accepted M tasks \\
    E10 & RQ2/4 & Unfiltered random $n{=}100$;
      \code{run\_e10\_m\_win\_v1} \\
    E12 & RQ4 & Seeds $\{0,1,2\}{\times}$120${\times}$B1/B2/M;
      \code{run\_e12\_m\_win\_v1} (1080 jobs) \\
    E14 & RQ1 & Length-matched feedback sweep;
      \code{run\_e14\_sweep\_v1} \\
    E16 & RQ4 & Second endpoint (\texttt{ecnu-max});
      \code{run\_e16\_canonical\_v1} \\
    E17 & RQ3 & Advisory gate \texttt{M\_adv} vs.\ A2 vs.\ M;
      \code{run\_e17\_advisory\_v1} \\
    \bottomrule
  \end{tabular}
\end{inlinetable}

%=============================================================================
\section{Statistical Protocol}
\label{sec:exp:stats}
%=============================================================================

Paired Wilcoxon signed-rank for conformance; Mann--Whitney~$U$ for latency;
Cliff's~$\delta$ for effect size; bootstrap 95\% CIs ($n{=}1000$);
$\alpha{=}0.05$.
Effect-size bands: $|\delta|{<}0.147$ negligible; ${<}0.33$ small;
${<}0.474$ medium; else large.
Holm--Bonferroni applies to \emph{pre-fix} multi-mode E1 comparisons on
\code{run\_e1\_canonical\_v1} only; fixed-oracle primary E1 reports means and
win rates without a Holm claim.

On the historical pre-fix corpus, observed M--B1/B2 Conf.\ gaps are
negligible and under-powered (post-hoc power ${\approx}3.6\%$ for
$|\delta|{=}0.014$; \texttt{power\_analysis.json}).
Fixed-oracle E1 emphasises per-task win rates; multi-seed
stability uses \code{run\_e12\_m\_win\_v1}, with E3/E10 stratification and E2
prevention as complementary evidence.

\subsection{Multi-Seed Protocol (E12)}
\label{sec:exp:e12}

E12-full re-runs 120 tasks ${\times}$ seeds $\{0,1,2\}$ ${\times}$ B1/B2/M under
the M-win protocol (\code{run\_e12\_m\_win\_v1}; M $K{=}5$).
Outcomes: per-seed mean Conf.\ ranking and pooled paired win rate M vs.\ B2.
B2 reaches 100\% Conf.\ on every seed; M averages 99.7\%---M does not lead
every seed, and B2 is slightly more seed-stable at the ceiling.

%=============================================================================
\section{Threats to Validity (Preview)}
\label{sec:exp:threats-preview}
%=============================================================================

Construct threats include finite Z3 witnesses (FAR bound in
Theorem~\ref{thm:far-bound}), incomplete pattern coverage, and first-order
mutants.
Internally, LLM stochasticity is mitigated by $T_\text{repair}{=}0.0$, fixed
seeds, bootstrap CIs, and E12; author-built benchmarks and the Z3 overlap
filter favour SMT-aided modes, addressed by E10, E8, and ablations.
Externally, primary evidence is synthetic numeric-domain Python with one LLM
family.
Conclusion validity separates pre-fix Holm results from fixed-oracle win
rates.
Full analysis: Chapter~\ref{ch:discussion}.

With the protocol fixed, Chapter~\ref{ch:results} reports mechanism first
(RQ1--RQ4), then quality--overhead and the E1 stress-test under RQ5.

% -----------------------------------------------------------------------------
\section*{Chapter Summary}
\addcontentsline{toc}{section}{Chapter Summary}
% -----------------------------------------------------------------------------

Five RQs and experiments E1--E17 test the Tier-1 claims on a 120-task hard
benchmark with ten compared modes.
Primary B1/B2/M ranking cites fixed-oracle M-win run ids
(Table~\ref{tab:corpus-card}); E6 and E2 use dedicated mechanism/prevention
corpora.
Metrics are Conf., strict success, PDR/FAR, and cost; statistics use
non-parametric tests with Holm correction only on the pre-fix multi-mode archive.
Threats are previewed here and fully analysed in Chapter~\ref{ch:discussion}.
